% ---
% id: meat-core/req:removal-checklist
% kind: req
% title: "Setting the removable flag"
% summary: >-
%   The engine flags a sector removable only when a witness it exhibits meets
%   all six conditions of the removal checklist, each decided exactly. Every
%   other outcome retains the sector, and no tolerance licenses a removal.
% prominence: standard
% status: implemented
% uses: [def:classes, thm:removal]
% mentions: [req:state-blind-test, rem:classes, rem:removal-general-format]
% assumes: []
% refs: []
% provenance: newly-derived
% implements: certiqs_key.meat_core.declaration.general_sectors:general_sector_report
% lean:
% ---

The flag rests on the verification of the checklist and on nothing about
where the witness came from, so the witness may come from any search. A
deletion on an approximate match would overstate the value
(Remark~\ref{rem:classes}), so every doubt retains the sector; and with a
\tref{party-local} entry or several
\tref[measured factor]{measured factors} the checklist has nothing to
certify or nothing to read (Remark~\ref{rem:removal-general-format}).

\begin{requirement}\label{req:removal-checklist}
The engine sets the flag \iface{removable} for a sector $s^{*}$ of a
declaration with one \tref{measured factor} when, and only when, a candidate
witness exhibited by the implementation is a witness for $s^{*}$ in the
sense of Definition~\ref{def:classes}\,(d): the list $\CList$ carries no
entry declared \tref{party-local} and the candidate meets all six of
(W1)--(W6); it sets the flag on no other ground. The
witness may come from any search, and it is never read as declared input: a
declaration supplies no witness. Each condition is decided exactly, at no
window: (W1) by the bit-level identity of
Requirement~\ref{req:state-blind-test}, (W2) as a comparison of sector
indices, (W5) as the finite combinatorial predicate of key-inertness, (W3)
and (W4) on the exhibited data as those conditions state, and (W6) on the
residual formed in the exhibited evaluation order; a witness whose
evaluation rounds anywhere, or whose residual is not the datum $0$, sets no
flag. Every other outcome --- a nonzero residual, a rounding operation, an
absent witness, a witness that names $s^{*}$, a target that is not
key-inert, a block that is not exactly scalar --- retains the sector. The
flag is set for no sector of a declaration whose list carries an entry
declared party-local, nor of a declaration with more than one measured
factor. No tolerance licenses a removal:
\iface{sector_criterion} is decided at exactly zero and not at the
\tref{declared tolerance} $\tnum$.
\end{requirement}
